\documentclass[10pt]{mypackage}

\usepackage[uprightgreeks,noDcommand]{kpfonts}
\renewcommand{\coloneq}{\coloneqq}
\renewcommand{\eqcolon}{\eqqcolon}
%\usepackage{homework}
\usepackage{notes}

\usepackage[ backend=bibtex, style = alphabetic, sorting=ynt ]{biblatex}
\addbibresource{exam_prep_references.bib}

\usepackage{parskip}

\fancyhf{}
\fancyhead[R]{Avinash Iyer}
\fancyhead[L]{Oral Exam Prep: Functional Analysis}
\fancyfoot[C]{\thepage}

\setcounter{secnumdepth}{0}

\begin{document}
\RaggedRight
\section{Locally Convex Topological Vector Spaces}%
\subsection{Definitions and Examples}%
\begin{definition}
  A topological vector space (TVS) is a vector space $X$ with a topology $ \mathcal{T} $ such that, with respect to this topology, the maps
  \begin{enumerate}[(i)]
    \item $X\times X \rightarrow X$, $\left( x,y \right)\mapsto x + y$
    \item $\F\times X \rightarrow X$, $\left( \alpha,x \right)\mapsto \alpha x$
  \end{enumerate}
  are continuous with respect to $ \mathcal{T} $.
\end{definition}
\begin{example}
  If $ \mathcal{P} $ is a family of seminorms on $X$, and $\mathcal{T}$ is the topology on $X$ with subbasic sets $U \left( x_0,p,\ve \right) = \set{x\in X : p\left( x-x_0 \right) < \ve}$, then $\left( X,\mathcal{T} \right)$ is a topological vector space.
\end{example}
\begin{definition}
  A locally convex topological vector space (LCS) is a TVS equipped with a family of seminorms $ \mathcal{P} $ that satisfy
  \begin{align*}
    \bigcap_{p\in \mathcal{P}} \set{x : p(x) = 0} = \set{0}.
  \end{align*}
  That is, the family $ \mathcal{P} $ separates the points of $X$.
\end{definition}
\begin{proposition}
  Let $X$ be a TVS, and let $ p $ be a seminorm on $X$. The following are equivalent:
  \begin{enumerate}[(a)]
    \item $p$ is continuous;
    \item $\set{x\in X : p(x) < 1}$  is open;
    \item $0\in \set{x\in X : p(x) < 1}^{\circ}$;
    \item $0\in \set{x\in X : p(x) \leq 1}^{\circ}$;
    \item $p$ is continuous at $0$;
    \item there is a continuous seminorm $q$ on $X$ such that $p\leq q$.
  \end{enumerate}
\end{proposition}
\begin{proof}
  We see immediately that (a) implies (b) since the set is equal to the inverse image of the open subset $[0,1)\subseteq [0,\infty)$. Similarly, if $p^{-1}\left( [0,1) \right)$ is open, then it is equal to its own interior, and since $0\in p^{-1}\left( [0,1) \right)$, it follows that there is a subbasic open neighborhood of $0$ contained in $p^{-1}\left( [0,1) \right)$, giving (c). Similarly, we obtain (d) by the fact that $0\in p^{-1}\left( [0,1) \right)^{\circ}\subseteq p^{-1}\left( [0,1] \right)^{\circ}$.

  Next, we see that (d) implies that for every $\ve > 0$, $0\in p^{-1}\left( [0,\ve] \right)^{\circ}$, so if $\left( x_i \right)_i$ is a net in $X$ converging to $0$ and $\ve > 0$, there is some $i_0$ such that for all $i\geq i_0$, $x_i\in \set{x : p(x) \leq \ve}$. That is, $p\left( x_i \right)\leq \ve$ for all $i\geq i_0$, whence $p$ is continuous at $0$. This gives (e).

  If $p$ is continuous at $0$, then if $x_i\rightarrow x$ is a convergent net in $X$, the reverse triangle inequality gives 
  \begin{align*}
    \left\vert p\left( x \right) - p\left( x_i \right) \right\vert &\leq p\left( x_i-x \right),
  \end{align*}
  whence $p\left( x_i \right) \rightarrow p\left( x \right)$, so $p$ is continuous. This gives that (e) implies (a).

  Taking $q = p$ yields (a) implies (f). We will show that (f) implies (e).If $x_i\rightarrow 0$ in $X$, then $q\left( x_i \right)\rightarrow 0$, but since $0\leq p\left( x_i \right)\leq q\left( x_i \right)$, it follows that $p\left( x_i \right)\rightarrow 0$.
\end{proof}
\begin{example}\hfill
  \begin{enumerate}[(a)]
    \item Let $X$ be a completely regular topological space, and let $C(X)$ be the set of all continuous functions from $X$ to $\F$. If $K$ is compact in $X$, we may define a seminorm $p_K\left( f \right) = \sup_{x\in K}\left\vert f(x) \right\vert$. Then, 
      \begin{align*}
        \set{p_K : K\subseteq X\text{ compact}}
      \end{align*}
      forms a separating family of seminorms for $C(X)$.
    \item Let $G$ be an open subset of $\C$. If $H(G)$ is the set of holomorphic functions on $G$, then the same seminorms $\set{p_K : K\subseteq G\text{ compact}}$ forms a locally compact space whose topology is uniform convergence on compact subsets.
    \item Let $X$ be a normed vector space. For each $\varphi\in X^{\ast}$, let $p_{\varphi}(x) = \left\vert \varphi(x) \right\vert$. Then, $ \mathcal{P} = \set{p_{\varphi} : \varphi\in X^{\ast}} $ forms a separating family of seminorms for $X$. The topology these seminorms induce is called the \textit{weak topology} on $X$. It is often denoted $\sigma\left( X,X^{\ast} \right)$.
    \item If $X$ is a normed space, then for each $x\in X$, we may define a seminorm $p_x\colon X^{\ast}\rightarrow [0,\infty)$ by taking $p_x\left( \varphi \right) = \left\vert \varphi(x) \right\vert$. This defines a separating family of seminorms $\mathcal{P} = \set{p_x : x\in X}$, and the topology on $X^{\ast}$ induced by these seminorms is called the weak*, or $w^{\ast}$, topology. It is often denoted $\sigma\left( X^{\ast},X \right)$.
  \end{enumerate}
\end{example}
Next, we focus on an examination of the ``convexity'' part of the notion of locally convex topological vector spaces.
\begin{definition}
  A subset $A\subseteq X$ is said to be convex if for all $v,w\in A$, the segment
  \begin{align*}
    \left[ v,w \right]\coloneq \set{tv + \left( 1-t \right)q : 0\leq t \leq 1}
  \end{align*}
  is contained in $A$.
\end{definition}
Note that the intersection of any family of convex subsets is convex.
\begin{definition}
  If $A\subseteq X$ is any subset, then the \textit{convex hull} of $A$, denoted $ \conv\left( A \right) $, is the intersection of all the convex subsets that contain $A$. If $X$ is a TVS, the \textit{closed convex hull} of $A$, denoted $ \overline{\conv}\left( A \right) $, is the intersection of all the closed convex sets that contain $A$.
\end{definition}
Given any linear map $T\colon X\rightarrow Y$ of real vector spaces, and $C\subseteq X$ a convex subset of $Y$, then $T^{-1}\left( C \right)$ is convex in $X$. If $X^{\ast}$ is the dual of any complex vector space, then since $X$ is also a real vector space, it then follows that the sets
\begin{align*}
  A_0 &= \set{x : \left\vert f(x) \right\vert \leq 1}\\
  A_1 &= \set{x : \re\left( f(x) \right) \leq 1}\\
  A_3 &= \set{x : \re\left( f(x) \right) > 1}
\end{align*}
are all convex.
\begin{proposition}
  Let $X$ be a TVS, and let $A$ be a convex subset of $X$. Then,
  \begin{enumerate}[(a)]
    \item $ \overline{A} $ is convex;
    \item if $a\in A^{\circ}$ and $b\in \overline{A}$, then
      \begin{align*}
        [a,b) \coloneq \set{tb + \left( 1-t \right)a : 0\leq t < 1}
      \end{align*}
      is contained in $A^{\circ}$.
  \end{enumerate}
\end{proposition}
\begin{proof}
  Let $a\in A$, $b\in \overline{A}$, and let $0\leq t \leq 1$. Let $\left( x_i \right)_i$ be a net in $A$ such that $x_i\rightarrow b$. Then, $tx_i + \left( 1-t \right)a \rightarrow tb + \left( 1-t \right)a$. In particular, this means that if $b\in \overline{A}$ and $a\in A$, then $\left[ a,b \right]\subseteq \overline{A}$. So, given any $a,b\in \overline{A}$, we have that the net $\left( x_i \right)_i\rightarrow b$ is such that $\left[x_i,a\right]\subseteq \overline{A}$, so it follows that $\left[ b,a \right]\subseteq \overline{A}$, whence $ \overline{A} $ is convex.

  For (b), fix $0 < t < 1$, and set $c = tb + \left( 1-t \right)a$ for some $a\in A^{\circ}$ and $b\in \overline{A}$. Since $a\in A^{\circ}$, there is an open neighborhood $V$ of $0$ such that $a + V\subseteq A$. Therefore, for any $d\in A$, we have
  \begin{align*}
    A &\supseteq td + \left( 1-t \right)\left( a + V \right)\\
      &= t\left( d-b \right) + tb + \left( 1-t \right)\left( a + V \right)\\
      &= \left( t\left( d-b \right) + \left( 1-t \right)V \right) + c.
  \end{align*}
  We hope to find some $d\in A$ such that $0\in t\left( d-b \right) + \left( 1-t \right)V$, which we call $U$. The preceding inclusion then shows that $c\in A^{\circ}$, since $U$ is an open neighborhood of $0$ with $c + U\subseteq A$.

  Finding such a $d$ is equivalent to finding $d$ such that $0\in \frac{1}{t}\left( 1-t \right)V + \left( d-b \right)$, or $d\in b-\frac{1}{t}\left( 1-t \right)V$,. Yet, we know that $0\in - \frac{1}{t}\left( 1-t \right)V$ by the assumption that $V$ is an open neighborhood of $0$. Since $b\in \overline{A}$, it follows that $b-\frac{1}{t}\left( 1-t \right)V \cap A \neq \emptyset$ by the definition of closure, whence such a $d\in A$ exists.
\end{proof}
\begin{corollary}
  If $A\subseteq X$, then $ \overline{\conv}(A) = \overline{\conv(A)} $.
\end{corollary}
\begin{definition}
  A set $A\subseteq X$ is called balanced if for any $ \left\vert \alpha \right\vert\leq 1 $ and any $x\in A$, we have $\alpha x\in A$.

  We say a set $A$ is \textit{absorbing} if for any $x\in X$, there is some $\ve > 0$ such that $tx\in A$ for all $0\leq t < \ve$. We say $A$ is absorbing at $a$ if for any $x\in X$, there is $\ve > 0$ such that $a + tx\in A$ for all $0\leq t < \ve$.
\end{definition}
Given a seminorm $p$ on $V$, the set $V = \set{x : p(x) < 1}$ is a convex balanced set that is absorbing at each of its points. We will show that the converse actually holds --- that is, a vector space is a LCS if and only if the set of all open, convex, balanced subsets of $X$ forms a neighborhood basis at $0$.
\begin{proposition}
  Suppose $X$ is a vector space over $F$, and $V$ is a nonempty convex balanced set that is absorbing at each of its points. Then, there is a unique seminorm $p$ on $X$ such that
  \begin{align*}
    V &= \set{x\in X : p(x) < 1}.
  \end{align*}
\end{proposition}
\begin{proof}
  Let
  \begin{align*}
    p(x) &= \inf \set{t\geq 0: x\in tV}.
  \end{align*}
  Since $V$ is absorbing, it follows that $X = \bigcup_{n=1}^{\infty}nV$, so this is a well-defined quantity. Furthermore, $p(0) = 0$.

  To show that $p\left( \alpha x \right) = \left\vert \alpha \right\vert p(x)$, we suppose $\alpha\neq 0$, and since $V$ is balanced,
  \begin{align*}
    p\left( \alpha x \right) &= \inf\set{t\geq 0 : \alpha x\in tV}\\
                             &= \inf\set{t \geq 0 : x\in \frac{t}{\alpha}V}\\
                             &= \inf\set{t\geq 0 : x\in \frac{t}{\left\vert \alpha \right\vert}V}\\
                             &= \left\vert \alpha \right\vert\inf\set{\frac{t}{\left\vert \alpha \right\vert} : x\in \frac{t}{\left\vert \alpha \right\vert}V}\\
                             &= \left\vert \alpha \right\vert p(x).
  \end{align*}
  We now observe that if $\alpha,\beta\geq 0$, then for any $a,b\in V$, we have
  \begin{align*}
    \alpha a + \beta b &= \left( \alpha + \beta \right) \left( \frac{\alpha}{\alpha + \beta}a + \frac{\beta}{\alpha + \beta}b \right)\\
                       &\in \left( \alpha + \beta \right)V,
  \end{align*}
  since $V$ is convex. For any $x,y\in X$, take $p(x) = \alpha$ and $p(y) = \beta$, and let $\delta > 0$. Then, $x\in \left( \alpha + \delta \right)V$ and $y\in \left( \beta + \delta \right)V$ by the definition of the infimum, so that $x + y\in \left( \alpha + \beta + 2\delta \right)V$. Therefore, we have $p\left( x + y \right)\leq \alpha + \beta + 2\delta$, which holds for all $\delta > 0$, so that $p\left( x + y \right)\leq p(x) + p(y)$.

  Finally, we will have to show that $V = \set{x : p(x) < 1}$. If $p(x) = \alpha < 1$, then for any $\alpha < \beta < 1$, we have $x\in \beta V\subseteq V$, which emerges from the fact that $V$ is balanced. Thus, $V \supseteq \set{x : p(x) < 1}$. If $x\in V$, then $p(x)\leq 1$, and since $V$ is absorbing at $x$, there is some $\ve > 0$ such that for $0 < t < \ve$, we have $y = x + tx\in V$. Yet, this means $x = \frac{1}{1+t}y$ with $y\in V$, so that $p(x) = \frac{1}{1 + t}p(y)\leq \frac{1}{1+t} < 1$.
\end{proof}
The seminorm $p$ is called the Minkowski function for $V$, or the \textit{gauge} of $V$.
\begin{proposition}
  Let $X$ be a TVS, and let $\mathcal{U}$ be the collection of open, convex, balanced subsets of $X$. Then, $X$ is locally convex if and only if $ \mathcal{U} $ is a basis for the neighborhood system at $0$.
\end{proposition}
\begin{proof}
  Suppose $X$ is locally convex with topology defined by the family of seminorms $ \mathcal{P} $. Then, we observe that each of the subbasic open neighborhoods of $0$, 
  \begin{align*}
    \set{U\left( 0,p,\ve \right) : \ve > 0,p\in \mathcal{P}},
  \end{align*}
  are open, convex, and balanced; since the set of all finite intersections of all these subbasic open neighborhoods yields a neighborhood basis for $0$, and finite intersections of open, balanced, convex subsets are open, balanced, and convex, so it follows that $\mathcal{U}$ contains the set of all finite intersections of all these subbasic open neighborhoods, whence $ \mathcal{U} $ contains a basis for the neighborhood system at $0$. In particular, this means $\mathcal{U}$ also must generate the neighborhood system at $0$, so $\mathcal{U}$ forms a basis for the neighborhood system at $0$.

  Now, suppose $\mathcal{U}$ forms a basis for the neighborhood system at $0$. Observe that an open, convex, balanced subset of $V$ contains $0$, and is itself absorbing. We claim that the family of Minkowski functionals
  \begin{align*}
    \mathcal{P} &= \set{p_V\left( \cdot - z \right) : z\in X, V\in \mathcal{U}}
  \end{align*}
  generates the topology for $X$.
\end{proof}
Thus, it follows that in a locally convex space, a net $\left( x_{i} \right)_i$ converges in the topology to $x$ if and only if for all seminorms $p\in \mathcal{P}$, we have $p\left( x_i - x \right)\rightarrow 0$.
\begin{example}\hfill
  \begin{enumerate}[(i)]
    \item Suppose $X$ is a locally compact Hausdorff space (or any Tychonoff space, but we will work with these for uniformity), and let $C(X)$ be the space of all continuous functions on $X$. Then, the family of seminorms $ p_K\left( f \right) = \sup_{x\in K} \left\vert f(x) \right\vert $ yields a locally convex space where convergence is defined as uniform convergence on compact subsets.
    \item If $X$ is a normed space, then the family of seminorms $p_{\varphi}\left( x \right) = \left\vert \varphi(x) \right\vert$ for each $\varphi\in X^{\ast}$ defines the weak topology on $X$.
    \item If $X$ is a normed space, then the family of seminorms $p_x(\varphi) = \left\vert \varphi(x) \right\vert$ for each $x\in X$ defines the weak*, or $w^{\ast}$, topology on $X^{\ast}$.
  \end{enumerate}
  The weak and weak* topologies are very important, and we will discuss them more in depth.
\end{example}
\subsection{Metrizability and Normability}%
When we look at a locally convex space, we concern ourselves with whether the topology of the space is given by a metric or not. It turns out this is (effectively) equivalent to there being countably many seminorms defining the topology.
\begin{proposition}
  Let $\set{p_n : n\geq 1}$ be a sequence of separating seminorms on $X$. For any $x,y\in X$, we may take
  \begin{align*}
    d\left( x,y \right) &= \sum_{n=1}^{\infty} 2^{-n} \frac{p_n\left( x-y \right)}{1 + p_n\left( x-y \right)}.
  \end{align*}
  Then, the topology on $X$ is defined by this metric.
\end{proposition}
\begin{definition}
  A Fréchet space is a locally convex space whose topology is defined by a metric that is complete.
\end{definition}
\begin{definition}
  If $X$ is a topological vector space, and $B\subseteq X$, then we say $B$ is bounded if, for every open set $U$ containing $0$, there is some $\ve > 0$ such that $\ve B \subseteq U$.
\end{definition}
\begin{proposition}
  If $X$ is a locally convex space, then $X$ is normable if and only if $X$ has a nonempty bounded open set.
\end{proposition}
\begin{proof}
  The open unit ball $U_X$ is bounded. Therefore, we only need to show one direction. Suppose $X$ is a locally convex space with a bounded open set $U$. Translating if necessary, we may assume $0\in U$.

  Since $X$ is locally convex, there is an open, convex, balanced subset of $U$ which we call $V$; using the gauge of $V$, we may find a continuous seminorm $p$ such that $\set{x : p(x) < 1} = V$. We will show that $p$ is a norm and that $p$ defines the topology on $X$.

  First, let $0\neq x\in X$. Let $W_0$ and $W_x$ be disjoint open subsets with $0\in W_0$ and $x\in W_x$. Since $W_0$ is absorbing, there is some $\ve > 0$ such that $W_0\supseteq \ve U \supseteq \ve V$. We note that $\ve V = \set{y : p(y) < \ve}$. Since $x\notin W_0$, we have $p(x)\geq \ve$, so that $p$ is a norm.

  To show that the topology on $X$ is generated by $p$, it suffices to show that if $q$ is any continuous seminorm, then there is some $\alpha$ such that $q\leq \alpha p$. Since $q$ is continuous, it follows that there is some $\ve > 0$ such that $\set{x : q(x) < 1}\supseteq \ve U \supseteq \ve V$. Yet, this means that if $p(x) < \ve$, then $q(x) < 1$, whence $q\leq \ve^{-1}p$.
\end{proof}
\subsection{Hyperplane Separation and Continuous Linear Functionals}%
We recall the classification of continuous linear functionals in a locally convex topological vector space.
\begin{theorem}
  If $X$ is a locally convex topological vector space with defining seminorms $ \mathcal{P} $, and $f\colon X\rightarrow \C$ is a linear functional, then the following statements are equivalent:
  \begin{enumerate}[(i)]
    \item $f$ is continuous;
    \item $f$ is continuous at $0$;
    \item $f$ is continuous at some point;
    \item $\ker(f)$ is closed;
    \item there are $p_1,\dots,p_n\in \mathcal{P}$ and some constant $C$ such that $\left\vert f(x) \right\vert\leq C\max\set{p_1(x),\dots,p_n(x)}$.
  \end{enumerate}
\end{theorem}
To understand hyperplane separation, we start with the most basic case --- namely, that of separating a convex subset from the origin via a continuous linear functional.
\begin{theorem}
  If $X$ is a topological vector space, and $G$ is an open convex nonempty subset of $X$ that does not contain the origin, then there is a closed hyperplane $M$ such that $M\cap G = \emptyset$.
\end{theorem}
\begin{proof}
  We start by assuming that $X$ is an $\R$-linear space. Choose $x_0\in G$, and let $H = x_0-G$. Then, $H$ is an open, convex subset of $X$ that contains $0$. Taking the gauge of $H$ yields a continuous sublinear functional $q\colon X\rightarrow \R$ with $H = \set{x : q(x) < 1}$. Furthermore, since $x_0\notin H$, we have that $q\left( x_0 \right)\geq 1$.

  Consider the subspace $Y = \set{\alpha x_0 : \alpha\in \R}$. Let $f_0\colon Y\rightarrow \R$ be defined by taking $f_0\left( \alpha x_0 \right) = \alpha q\left( x_0 \right)$. Observe that if $\alpha \geq 0$, then $f_0\left( \alpha x_0 \right) = q\left( \alpha x_0 \right)$, and if $\alpha < 0$, then $f_0\left( \alpha x_0 \right) = \alpha q\left( x_0 \right) \leq \alpha < 0 \leq q\left( \alpha x_0 \right)$, whence $f_0\leq q$ on $Y$. Using the Hahn--Banach--Minkowski extension yields some $f\colon X\rightarrow \R $ a continuous linear functional such that $f|_{Y} = f_0$ and $f\leq q$ on $X$. Then, $M = \ker\left( f \right)$ is the desired hyperplane.

  If $x\in G$, then $x_0-x\in H$, so $f\left( x_0 \right) - f\left( x \right) = f\left( x_0-x \right) \leq q\left( x_0-x \right) < 1$, while $f(x) > f\left( x_0 \right) - 1 = q\left( x_0 \right)-1\geq 0$ for all $x\in G$, whence $M \cap G = \emptyset$.

  If $X$ is a $\C$-linear space, then since $X$ is also an $\R$-linear space, we obtain a continuous $\R$-linear functional $f\colon X\rightarrow \R$ with $G\cap \ker\left( f \right) = \emptyset$. If $F(x) = f(x) - if(ix)$, then $F$ is a $\C$-linear functional with $f = \re(F)$ with $F(x) = 0$ if and only if $f(x) = f(ix) = 0$, so that $M = \ker\left( F \right)$ is the desired closed hyperplane.
\end{proof}
Affine hyperplanes in $X$ are given by $M' = M + x_0$ for some $x_0\in X$ and hyperplane $M$. It turns out that we can extend any separating affine (closed) subspace of a topological vector space to an affine hyperplane.
\begin{corollary}
  Let $X$ be a topological vector space, and let $G$ be an open, convex, nonempty subset of $X$. If $Y$ is an affine subspace of $X$ with $Y\cap G = \emptyset$, then there is a closed affine hyperplane $M$ in $X$ such that $Y\subseteq M$ and $M\cap G = \emptyset$.
\end{corollary}
\begin{proof}
  Choosing some $x_0\in Y$, we may consider $G-x_0$ and $Y - x_0$, so that $Y$ may be assumed to be a subspace of $X$. Let $Q\colon X\rightarrow X/Y$ be the quotient map. Notice that $Q^{-1}\left( Q(G) \right) = \set{y + G : y\in Y}$, so that $Q(G)$ is open in $X/Y$. Furthermore, $Q(G)$ is also convex since $Q$ is a linear map. Furthermore, since $Y\cap G = \emptyset$, $0\notin Q(G)$. In particular, there is then a closed hyperplane $N$ in $X/Y$ such that $N\cap Q(G) = \emptyset$. Setting $M = Q^{-1}\left( N \right)$, we have that $M$ is our desired hyperplane.
\end{proof}
In the special case of a real topological vector space, it's useful to note that $X\setminus \ker\left( f \right)$ for a continuous linear functional $f\colon X\rightarrow \R$ has two connected components. This yields the following notion.
\begin{definition}
  Let $X$ be a real topological vector space. A subset $S\subseteq X$ is called an open half space if there is some continuous linear functional $f\colon X\rightarrow \R$ such that $S = f^{-1}\left( \left( \alpha,\infty \right) \right)$. We call $S$ a closed half-space if there is a continuous linear functional $f\colon X\rightarrow \R$ such that $S = f^{-1}\left( [\alpha,\infty) \right)$.
\end{definition}
We say $A$ and $B$ are strictly separated if they are contained in disjoint open half-spaces, while they are separated if they are contained in two closed half-spaces whose intersection is of the form $M = f^{-1}\left( \set{\alpha} \right)$. That is, their intersection is a closed affine hyperplane. Notice also that continuous linear functionals are open maps, so that the interior of a closed half-space is an open half-space, and the closure of an open half-space is a closed half-space.
\begin{theorem}[General Hyperplane Separation]
  Let $X$ be a real topological vector space, and $A$ and $B$ disjoint convex subsets with $A$ open. Then, there is a continuous linear functional $f\colon X\rightarrow \R$ and a real scalar $\alpha$ such that $A\subseteq f^{-1}\left( \left( -\infty,\alpha \right) \right)$ and $B\subseteq f^{-1}\left( [\alpha,\infty) \right)$.
\end{theorem}
\begin{proof}
  Let $G = A - B$. Then, $G$ is convex, and $G = \bigcup_{b\in B}\left( A-b \right)$, meaning $G$ is open. Since $A\cap B = \emptyset$, it follows that $0\notin G$. Thus, there is a closed hyperplane $M$ in $X$ such that $M\cap G = \emptyset$.

  Let $f\colon X\rightarrow \R$ be a linear functional such that $M = \ker\left( f \right)$. Notice that $f(G)$ is a convex subset of $\R$ with $0\notin f(G)$. Either, $f(x) > 0$ for all $x\in G $ or $f(x) < 0$ for all $x\in G$. Without loss of generality, we may assume that $f(x) > 0$ for all $x\in G$, so that $f(b) < f(a)$ for any $a\in A$ and $b\in B$, meaning there is some real number $\alpha$ such that $\sup_{b\in B}f(b)\leq \alpha \leq \inf_{a\in A}f(a)$. Since $f(A)$ and $f(B)$ are open intervals if $B$ is open, it follows that $f < \alpha$ on $B$ and $f > \alpha$ on $A$.
\end{proof}
\begin{lemma}
  If $X$ is a topological vector space, $K$ is a compact subset of $X$, and $V$ is an open subset satisfying $K\subseteq V$, then there is an open neighborhood $U$ of $0$ such that $K + U \subseteq V$.
\end{lemma}
\begin{proof}
  Let $ \mathcal{O}_0 $  denote the system of open neighborhoods of $0$. Suppose toward contradiction that for each $U\in \mathcal{O}_0$ that $K + U$ is not contained in $V$. We may choose vectors $x_U\in K$ and $y_U\in U$ such that $x_U + y_U\in V^{c}$. This yields a net $\left( y_U \right)_{U\in \mathcal{O}_0}$ and $\left( x_U \right)_{U\in \mathcal{O}_0}$ with $y_U\rightarrow 0$ and a subnet $x_{U'}\rightarrow x\in K$. Yet, this yields $\left( x_{U'} + y_{U'} \right)\rightarrow x$, implying that $x\in \overline{V^c} = V^c$, which yields a contradiction.
\end{proof}
\begin{theorem}[Hyperplane Separation for Locally Convex Spaces]
  Let $X$ be a real locally convex topological vector space, and let $A$ and $B$ be two disjoint closed convex subsets of $X$ with $B$ compact. Then, $A$ and $B$ are contained in disjoint open half-spaces.
\end{theorem}
\begin{proof}
  Notice that $B$ is a compact subset of $A^{c}$. Therefore, there is an open neighborhood $U_1$ of $0$ such that $B + U_1\subseteq A^{c}$. Since $X$ is locally convex, we may find a continuous seminorm $p$ on $X$ such that $\set{x : p(x) < 1}\subseteq U_1$.

  Define $U = \set{x : p(x) < 1/2}$. Then, $\left( B + U \right) \cap \left( A + U \right) = \emptyset$, and $A + U$ and $B + U$ are open convex subsets of $X$ that contain $A$ and $B$. Therefore, the result follows from the general hyperplane separation.
\end{proof}
Rephrasing this theorem for locally convex spaces by using the fact that any linear functional $\varphi$ over $\C$ is given by $\varphi(x) = \re\left( \varphi(x) \right) - i\re\left( \varphi\left( ix \right) \right)$ yields the following.
\begin{theorem}[Hyperplane Separation for Locally Convex Spaces, Encore]
  Let $X$ be a locally convex space, and let $A$ and $B$ be disjoint closed convex subsets of $X$ with $B$ compact. Then, there is a continuous linear functional $f\colon X\rightarrow \C$, some $\alpha\in \R$, and some $\ve > 0$ such that for all $a\in A$ and $b\in B$, we have
  \begin{align*}
    \re\left( f(a) \right)\leq \alpha < \alpha + \ve \leq \re\left( f(b) \right).
  \end{align*}
\end{theorem}
We will now discuss the dual space of one particular locally convex space.
\begin{proposition}
  Let $X$ be a locally compact Hausdorff space, and let $C(X)$ be topologized by uniform convergence on compact subsets. If $L\colon C(X)\rightarrow \F$ is a continuous linear functional, then there is a compact set $K$ and a regular (complex) Borel measure $\mu$ on $K$ such that $L(f) = \int_{K}^{} f\:d\mu$ for any $f\in C(X)$. Conversely, all these measures define elements of $C(X)^{\ast}$.
\end{proposition}
\begin{proof}
  Any compactly support measure defines an element of $C(X)^{\ast}$, and satisfies $\left\vert L(f) \right\vert \leq \norm{\mu}\sup_{x\in K}\left\vert f(x) \right\vert$ for any such measure $\mu$.

  Conversely, let $L\in C(X)^{\ast}$. By the classification of continuous linear functionals, there are compact sets $K_1,\dots,K_n$ and a constant $C$ such that $\left\vert L(f) \right\vert\leq C \max\set{p_{K_j}(f) : j=1,\dots,n}$. Let $K = \bigcup_{j=1}^{n}K_j$. Then, $\left\vert L(f) \right\vert\leq \alpha p_K(f)$. In particular, for any $f\in C(X)$ with $f|_{K} = 0$, then $L(f) = 0$.

  Define a linear functional $F\colon C(K)\rightarrow \F$ as follows. For any $g\in C(K)$, let $ \widetilde{g} $ be a continuous extension to $X$, and set $F(g) = L\left( \widetilde{g} \right)$. Suppose $ \widetilde{g}_1 $ and $ \widetilde{g}_2 $ are any such extensions; then, $g_1-g_2 = 0$ on $K$, so that $L\left( \widetilde{g}_1 \right) = L\left( \widetilde{g}_2 \right)$, so that $F$ is well-defined. Furthermore, $F$ is linear. Since, for any $g\in C(K)$ and extension $ \widetilde{g}\in C(X) $, we have $\left\vert F(g) \right\vert \leq C p_{K}(g) = C\norm{g}$. In particular, this means $F$ is continuous, so there is some regular complex Borel measure $\mu$ in $M(K)$ such that $F(g) = \int_{}^{} g\:d\mu$. Notice that if $f\in C(X)$ is such that $f|_{K} = g\in C(K)$, it follows that $L(f) = F(g) = \int_{}^{} f\:d\mu$.
\end{proof}
\section{Weak Topologies and Extremal Structure}%
\subsection{Duality}%
Let $X$ be a locally convex space. Then, $X^{\ast}$ denotes the space of continuous linear functionals on $X$, and admits a vector space structure. We will write $\left\langle x,\varphi \right\rangle = \varphi(x)$ whenever $x\in X$ and $\varphi\in X^{\ast}$. We will write $ \hat{x} = \left\langle x,\cdot \right\rangle $ to denote the analogous linear functional on $\varphi$ such that $\hat{x}(\varphi) = \left\langle x,\varphi \right\rangle$.
\begin{definition}
  If $X$ is a locally convex space, then the weak topology on $X$ is the topology defined by the family of seminorms
  \begin{align*}
    p_{\varphi}(x) &= \left\vert \left\langle x,\varphi \right\rangle \right\vert.
  \end{align*}
  The weak* topology on $X$ is the topology defined by the seminorms
  \begin{align*}
    p_{x}\left( \varphi \right) &= \left\vert \left\langle x,\varphi \right\rangle \right\vert.
  \end{align*}
  We have that $\left( x_i \right)_{i}\xrightarrow{\operatorname{wk}} x$ if, for all $\varphi\in X^{\ast}$, $\varphi\left( x_i \right)\rightarrow \varphi(x)$. Similarly, we have that $\left( \varphi_i \right)_i\xrightarrow{\operatorname{wk}^{\ast}} \varphi$ if, for all $x\in X$, $\varphi_i(x)\rightarrow \varphi(x)$.
\end{definition}
\begin{theorem}
  If $X$ is a locally convex space, then $\left( X^{\ast},\operatorname{wk}^{\ast} \right)^{\ast} = X$.
\end{theorem}
\begin{proof}
  First, if $x\in X$, then $ \varphi\mapsto \left\langle x,\varphi \right\rangle $ is a weak*-continuous linear functional on $X^{\ast}$ by definition.

  Conversely, if $f\in \left( X^{\ast},\operatorname{wk}^{\ast} \right)^{\ast}$, then there are vectors $x_1,\dots,x_n\in X$ such that
  \begin{align*}
    \left\vert f\left( \varphi \right) \right\vert &\leq \sum_{k=1}^{n} \left\vert \left\langle x_k,\varphi \right\rangle \right\vert
  \end{align*}
  for all $\varphi\in X^{\ast}$. This implies that
  \begin{align*}
    \bigcap_{k=1}^{n} \ker\left( \hat{x}_{k} \right) &\subseteq \ker\left( f \right),
  \end{align*}
  so there are $\alpha_1,\dots,\alpha_n$ such that $f = \sum_{k=1}^{n}\alpha_k\hat{x}_k$, meaning that $f\in X$. 
\end{proof}
Local convexity is a very useful property to have, as the following result exhibits.
\begin{theorem}
  Let $X$ be a locally convex space, and let $A$ be a convex subset of $X$. Then, $ \overline{A} = \overline{A}^{\operatorname{wk}} $.
\end{theorem}
\begin{proof}
  Notice that the weak topology is weaker than the topology on $X$, so $ \overline{A}\subseteq \overline{A}^{\operatorname{wk}} $.

  Conversely, suppose there were $x\in \overline{A}^{\operatorname{wk}}\setminus \overline{A}$. Then, by geometric Hahn--Banach, there is $\varphi\in X^{\ast}$, $\alpha\in \R$, and $\ve > 0$ such that
  \begin{align*}
    \re\left( \left\langle a,\varphi \right\rangle \right) \leq \alpha < \alpha + \ve \leq \re\left( \iprod{x}{\varphi} \right)
  \end{align*}
  for all $a\in \overline{A}$. Yet, this means that $ \overline{A}\subseteq \set{y\in X : \re\left( \left\langle y,\varphi \right\rangle \right)\leq \alpha} $, even though the latter is weakly closed. Since $x\notin B$, it follows that $x\notin \overline{A}$, which yields a contradiction.
\end{proof}
\subsection{Duality in Subspaces and Quotient Spaces}%
\begin{proposition}
  Let $p$ be a seminorm on $ X $, $M$ a subspace of $X$, and $ \overline{p}\colon X/M\rightarrow [0,\infty) $ is defined by
  \begin{align*}
    \overline{p}\left( x + M \right) &= \inf\set{p\left( x + y \right) : y\in M},
  \end{align*}
  then $ \overline{p} $ is a seminorm on $X/M$. If $ \mathcal{P} $ is a family of continuous seminorms on $X$ defining a locally convex topology, then $ \overline{ \mathcal{P} } $ defines the quotient topology on $X/M$.
\end{proposition}
\section{$C^{\ast}$-Algebras}%
\section{Spectral Theory}%
Spectral theory is primarily concerned with representations of abelian $C^{\ast}$-algebras. Consider a normal operator $N$ on $H$. Then, $C^{\ast}\left( N \right)$ is an abelian $C^{\ast}$-algebra, and the continuous functional calculus $f\mapsto f\left( N \right)$ yields a $\ast$-isomorphism of $C\left( \sigma(N) \right)$ onto $C^{\ast}\left( N \right)$. In particular, $f\mapsto f\left( N \right)$ is a representation of the abelian $C^{\ast}$-algebra $C^{\ast}\left( N \right)$.

We will find a way to extend this further, as well as find some useful applications to the structure of abelian von Neumann algebras and understanding unitary invariants for normal operators.
\begin{definition}
  Let $X$ be a set, $\Sigma$ a $\sigma$-algebra on $X$, and $H$ a Hilbert space. A \textit{spectral measure} for $\left( X,\Sigma,H \right)$ is a map $E\colon \Sigma\rightarrow B\left( H \right)$ satisfying
  \begin{enumerate}[(a)]
    \item for each $\Delta\in \Sigma$, $E\left( \Delta \right)$ is a projection;
    \item $E\left( \emptyset \right) = 0$ and $E\left( X \right) = 1$;
    \item $E\left( \Delta_1\cap \Delta_2 \right) = E\left( \Delta_1 \right)E\left( \Delta_2 \right)$ for any $\Delta_1,\Delta_2\in \Sigma$;
    \item given pairwise disjoint $\set{\Delta_n}_{n=1}^{\infty}$, we have
      \begin{align*}
        E\left( \bigcup_{n=1}^{\infty}\Delta_n \right) &= \sum_{n=1}^{\infty} E\left( \Delta_n \right),
      \end{align*}
      where the latter is an SOT-convergent sum of pairwise orthogonal projections converging to the projection onto the supremum $\bigwedge_{n=1}^{\infty}E\left( \Delta_n \right)H$.
  \end{enumerate}
\end{definition}
Recall that the SOT and WOT are defined by the seminorms
\begin{itemize}
  \item $p_h\left( A \right) = \norm{Ah}$ for the SOT, and
  \item $p_{h,k}\left( H \right) = \left\vert \iprod{Ah}{k} \right\vert$ for the WOT.
\end{itemize}
\begin{example}
  If $\left( X,\Omega,\mu \right)$ is a $\sigma$-finite measure space, and $\phi\in L^{\infty}\left( X,\mu \right)$, then $M_{\phi}\in B\left( L^2\left( X,\mu \right) \right)$ can be defined by $M_{\phi}\left( f \right) = \phi f$.

  A net $\left( \phi_i \right)_{i\in I}$ in $L^{\infty}\left( X,\mu \right)$ is $w^{\ast}$-convergent to $\phi$ if and only if $M_{\phi_i}\rightarrow M_{\phi}$ in WOT. This follows from the fact that if $f,g\in L^2\left( X,\mu \right)$, then $f \overline{g}\in L^1\left( X,\mu \right)$, where we have $L^1\left( X,\mu \right)^{\ast} = L^{\infty}\left( X,\mu \right)$, so that if $\left( \phi_i \right)_i\rightarrow \phi$ weak*, then
  \begin{align*}
    \iprod{M_{\phi_i}f}{g} &= \int_{}^{} \phi_i \left( f \overline{g} \right)\:d\mu\\
                           &\rightarrow \int_{}^{} \phi f \overline{g}\:d\mu\\
                           &= \iprod{M_{\phi}f}{g},
  \end{align*}
  while the converse emerges from the fact that any element $f\in L^1\left( X,\mu \right)$ can be written as $ f = g_1 \overline{g_2} $.
\end{example}
\begin{example}
  If $E$ is a spectral measure for $\left( X,\Omega,H \right)$, then $E^{(n)}\left( \Delta \right)= E\left( \Delta \right)^{(n)}$ (under the diagonal embedding) is a spectral measure for $\left( X,\Omega,H^{(n)} \right)$.
\end{example}
We will now use spectral measures to define a representation of a $C^{\ast}$-algebra.
\begin{lemma}
  Let $E$ be a spectral measure on $\left( X,\Omega,H \right)$, and let $g,h\in H$. Then,
  \begin{align*}
    E_{g,h} \left( \Delta \right) &= \iprod{E\left( \Delta \right)g}{h}
  \end{align*}
  defines a countably additive measure on $\Omega$ with total variation bounded above by $\norm{g}\norm{h}$.
\end{lemma}
\begin{proof}
  For countable additivity, we note the classic result that a sequence of projections $\left( P_n \right)_n\rightarrow P$ converges in WOT if and only if it converges in SOT. In particular, this means that if $\set{\Delta_n}_{n\in \N}$ is a sequence of disjoint subsets of $X$, then we have a sequence of projections
  \begin{align*}
    \left( \sum_{k=0}^{n}E\left( \Delta_k \right) \right)_{n} &\rightarrow \sum_{k=0}^{\infty} E\left( \Delta_k \right)
  \end{align*}
  converging in SOT, hence converging in WOT. Yet, this means that
  \begin{align*}
    \iprod{\left( \sum_{k=0}^{n}E\left( \Delta_k \right) \right)g}{h} &= \sum_{k=0}^{n} \iprod{E\left( \Delta_k \right)g}{h}\\
                                                                      &\rightarrow  \iprod{\left( \sum_{k=0}^{\infty}E\left( \Delta_k \right) \right)g}{h},
  \end{align*}
  so that
  \begin{align*}
    \sum_{k=0}^{\infty} \iprod{E\left( \Delta_k \right)g}{h} &= \iprod{E\left( \bigcup_{k=0}^{\infty}\Delta_k \right)g}{h},
  \end{align*}
  yielding our desired result.

  If we have a finite set $\Delta_1,\dots,\Delta_n$ of pairwise disjoint sets, we choose $\left\vert \alpha_j \right\vert = 1$ such that $ \left\vert \iprod{E\left( \Delta_j \right)g}{h} \right\vert = \alpha_j \iprod{E\left( \Delta_j \right)g}{h}$, whence
  \begin{align*}
    \sum_{j=1}^{n} \left\vert \mu\left( \Delta_j \right) \right\vert &= \iprod{\sum_{j=1}^{n}\alpha_j E\left( \Delta_j \right)g}{h}\\
                                                                     &\leq \norm{\sum_{j=1}^{n}E\left( \Delta_j \right)\alpha_j g}\norm{h}.
  \end{align*}
  Notice that the family of $E\left( \Delta_j \right)$ are pairwise orthogonal projections, meaning that by Pythagoras's Theorem and the additivity of the spectral measure, we have
  \begin{align*}
    \norm{\sum_{j}^{n}E\left( \Delta_j \right)\alpha_jg}^2 &= \sum_{j=1}^{n}\norm{E\left( \Delta_j \right)g}^2\\
                                                           &= \norm{E\left( \bigcup_{j=1}^{n}\Delta_j \right)g}^2\\
                                                           &\leq \norm{g}^2,
  \end{align*}
  whence $\sum_{j=1}^{n}\left\vert \mu\left( \Delta_j \right) \right\vert\leq \norm{g}\norm{h}$. Since this holds for any finite collection of subsets, we have $\norm{\mu}_{\operatorname{TV}}\leq \norm{g}\norm{h}$.
\end{proof}
We next concern ourselves with establishing ``integration'' with respect to the spectral measure.
\begin{proposition}
  Let $E$ be a spectral measure for $\left( X,\Omega,H \right)$, and let $\phi\colon X\rightarrow \C$ be a bounded $\Omega$-measurable function. Then, there is a unique operator $A\in B(H)$ such that for any $\ve > 0$ and partition $\set{\Delta_1,\dots,\Delta_n}$ with $\sup_{x,x'\in \Delta_k}\left\vert \phi(x)-\phi(x') \right\vert < \ve$ for each $k$, we have
  \begin{align*}
    \norm{A - \sum_{k=1}^{n}\phi\left( x_k \right)E\left( \Delta_k \right)} < \ve
  \end{align*}
  for any $x_k\in \Delta_k$.
\end{proposition}
\begin{proof}
  For each $g,h\in H$, define a (bounded) sesquilinear form
  \begin{align*}
    B\left( g,h \right) &= \int_{}^{} \phi\:dE_{g,h},
  \end{align*}
  with $\left\vert B\left( g,h \right) \right\vert \leq \norm{\phi}_{L^{\infty}}\norm{g}\norm{h}$. Then, by the structure of bounded sesquilinear forms, there is a unique operator $A$ such that $ B\left( g,h \right) = \iprod{Ag}{h} $ for all $g,h\in H$.

  Consider a partition $\set{\Delta_1,\dots,\Delta_n}$ satisfying the conditions in the proposition. For any arbitrary vectors $g,h\in H$ and $x_k\in \Delta_k$, we have
  \begin{align*}
    \left\vert \iprod{Ag}{h} - \sum_{k=1}^{n} \phi\left( x_k \right) \iprod{E\left( \Delta_k \right)g}{h} \right\vert &= \left\vert \sum_{k=1}^{n} \int_{\Delta_k}^{} \left( \phi(x)-\phi\left( x_k \right) \right)\:dE_{g,h} \right\vert\\
                  &\leq \sum_{k=1}^{n} \int_{\Delta_k}^{} \left\vert \phi(x)-\phi\left( x_k \right) \right\vert\:d\left\vert E_{g,h} \right\vert\\
                  &\leq \ve \int_{}^{} \:d\left\vert E_{g,h} \right\vert\\
                  &\leq \ve \norm{g}\norm{h}.
  \end{align*}
  This yields the desired result.
\end{proof}
For the operator $A$ in the above statement, we write
\begin{align*}
  A &= \int_{}^{} \phi\:dE.
\end{align*}
That is, for any $g,h\in H$ and bounded $\Omega$-measurable function $\phi$ on $X$, we have
\begin{align*}
  \iprod{\left( \int_{}^{} \phi\:dE \right)g}{h} &= \int_{}^{} \phi\:dE_{g,h}.
\end{align*}
If $B_{\infty}\left( X,\Omega \right)$ denotes the set of bounded $\Omega$-measurable functions $\phi\colon X\rightarrow \C$ with the norm $\norm{\phi} = \sup_{x\in X} \left\vert \phi(x) \right\vert$, then $B_{\infty}\left( X,\Omega \right)$ is a unital Banach algebra --- in fact, it is an abelian $C^{\ast}$-algebra.
\begin{proposition}
  Suppose $E$ is a spectral measure for $\left( X,\Omega, H \right)$, and $\rho\colon B_{\infty}\left( X,\Omega \right)$ is defined by taking $\rho(\phi) = \int_{}^{} \phi\:dE$.

  Then, $\rho$ is a representation of $B_{\infty}\left( X,\Omega \right)$ and $\rho(\phi)$ is a normal operator for every $\phi\in B_{\infty}\left( X,\Omega \right)$.
\end{proposition}
\begin{proof}
  We will only focus on showing that $\rho$ is multiplicative. Suppose $\phi,\psi\in B_{\infty}\left( X,\Omega \right)$. Let $\ve > 0$, and let $\set{\Delta_1,\dots,\Delta_n}$ be a Borel partition of $X$ satisfying
  \begin{align*}
    \sup\set{\left\vert \omega(x)-\omega(x') \right\vert : x,x'\in\Delta_k,\omega=\phi,\psi,\phi\psi} &< \ve.
  \end{align*}
  In particular, we have that for any $x_k\in \Delta_k$ and $\omega = \phi,\psi,\phi\psi$, we have
  \begin{align*}
    \norm{ \int_{}^{} \omega\:dE - \sum_{k=1}^{n} \omega\left( x_k \right)E\left( \Delta_k \right) } &< \ve.
  \end{align*}
  The triangle inequality then gives
  \begin{align*}
    \norm{ \int_{}^{} \phi\psi\:dE - \left( \int_{}^{} \phi\:dE \right) \left( \int_{}^{} \psi\:dE \right) } &\leq \ve + \norm{\sum_{k=1}^{n}\phi\left( x_k \right)\psi\left( x_k \right)E\left( \Delta_k \right) - \left( \sum_{i=1}^{n}\phi\left( x_i \right)E\left( \Delta_i \right) \right)\left( \sum_{j=1}^{n}\psi\left( x_j \right)E\left( \Delta_{x_j} \right) \right)}\\
                & + \norm{ \left( \sum_{i=1}^{n}\phi\left( x_i \right)E\left( \Delta_i \right) \right)\left( \sum_{j=1}^{n} \psi\left( x_j \right)E\left( \Delta_j \right) \right) - \left( \int_{}^{} \phi\:dE \right)\left( \int_{}^{} \psi\:dE \right) }.
  \end{align*}
  Since $E\left( \Delta_i \right)E\left( \Delta_j \right) = E\left( \Delta_i\cap \Delta_j \right)$, and $\set{\Delta_1,\dots,\Delta_n}$ is a partition, we have that the second term on the right-hand side is zero. Thus, we have
  \begin{align*}
    \norm{ \int_{}^{} \phi\psi\:dE - \left( \int_{}^{} \phi\:dE \right) \left( \int_{}^{} \psi\:dE \right) } &\leq \ve + \norm{ \left( \sum_{i=1}^{n}\phi\left( x_i \right)E\left( \Delta_i \right) \right) \left( \sum_{j=1}^{n} \psi\left( x_j \right)E\left( \Delta_j \right) - \int_{}^{} \psi\:dE \right) }\\
                &+ \norm{ \left( \sum_{i=1}^{n}\phi\left( x_i \right)E\left( \Delta_i \right) - \int_{}^{} \phi\:dE\right) \left( \int_{}^{} \psi\:dE \right) }\\
                &\leq \ve\left( 1 + \norm{\phi} + \norm{\psi} \right).
  \end{align*}
  Thus, we get multiplicativity.
\end{proof}
\begin{corollary}
  If $X$ is a compact Hausdorff space, and $E$ is a spectral measure on $H$ defined on the Borel sets of $X$, then $\rho\colon C(X)\rightarrow B(H)$ given by $ \rho(u) = \int_{}^{} u\:dE $ is a representation of $C(X)$.
\end{corollary}
Thus far, all we have obtained is a variety of the properties of a spectral measure. However, we do not know if such a measure exists. The following is the proof that it does.
\begin{theorem}
  Suppose $\rho\colon C(X)\rightarrow B(H)$ is a representation. Then, there is a unique spectral measure $E$ defined on the Borel subsets of $X$ such that for all $g,h\in H$, $E_{g,h}$ is a regular measure, and we have
  \begin{align*}
    \rho(u) &= \int_{}^{} u\:dE
  \end{align*}
  for all $u\in C(X)$.
\end{theorem}
\begin{theorem}
  Let $g,h\in H$. The map $u\mapsto \iprod{\rho(u)g}{h}$ is a linear functional on $C(X)$ with norm at most $\norm{g}\norm{h}$. So, there is a unique measure $\mu_{g,h}$ in $M(X)$ such that
  \begin{align*}
    \iprod{\rho(u)g}{h} &= \int_{}^{} u\:d\mu_{g,h}.
  \end{align*}
  The pairing $ \left( g,h \right)\mapsto \mu_{g,h} $ is sesquilinear, and for fixed $\phi\in B_{\infty}(X)$, the form
  \begin{align*}
    \left[ g,h \right] &= \int_{}^{} \phi\:d\mu_{g,h}
  \end{align*}
  is a sesquilinear form with $ \norm{\left[ g,h \right]}\leq \norm{\phi}\norm{g}\norm{h} $. Therefore, there is a unique bounded operator $A$ such that $ \left[ g,h \right] = \iprod{Ag}{h} $ with $\norm{A}\leq \norm{\phi}$. We will denote this operator $A$ by $ \widetilde{\rho}(\phi) $. We obtain a well-defined map $ \widetilde{\rho}\colon B_{\infty}(X)\rightarrow B(H) $ satisfying $ \norm{ \widetilde{\rho}(\phi) }\leq \norm{\phi} $ such that
  \begin{align*}
    \iprod{ \widetilde{\rho}(\phi)g }{h} &= \int_{}^{} \phi\:d\mu_{g,h}.
  \end{align*}
  We start by claiming that $ \widetilde{\rho}\colon B_{\infty}(X)\rightarrow B(H) $ is a representation with $ \widetilde{\rho}|_{C(X)} = \rho $. That $ \widetilde{\rho} $ restricted to $C(X)$ is equal to $\rho$ follows directly from the definition of $ \widetilde{\rho} $.

  Now, let $\phi\in B_{\infty}(X)$. Consider $ \phi $ as an element of $ M(X)^{\ast} = C(X)^{\ast\ast}$ by viewing it as corresponding to the linear functional
  \begin{align*}
    \mu &\mapsto \int_{}^{} \phi\:d\mu.
  \end{align*}
  By the Banach--Alaoglu theorem, we have that $\set{u\in C(X) : \norm{u}\leq \norm{\phi}}$ is weak*-dense in $\set{L\in M(X)^{\ast} : \norm{L}\leq \norm{\phi}}$. Then, there is a net $\left( u_i \right)_i\subseteq C(X)$ with $\norm{u_i}\leq \norm{\phi}$ for all $i$ and
  \begin{align*}
    \int_{}^{} \phi\:d\mu &= \lim_{i} \int_{}^{} u_i\:d\mu
  \end{align*}
  for all $\mu\in M(X)$. For any $\psi\in B_{\infty}(X)$, we have $\psi\mu \in M(X)$ for any $\mu\in M(X)$. In particular, we have
  \begin{align*}
    \int_{}^{} \phi\psi\:d\mu &= \lim_{i} \int_{}^{} u_i\psi\:d\mu
  \end{align*}
  for all $\psi\in B_{\infty}(X)$ and $\mu\in M(X)$, whence $ \widetilde{\rho}\left( u_i\psi \right)\rightarrow \widetilde{\rho}\left( \phi\psi \right) $ in WOT. For any $\psi\in C(X)$, we have that 
  \begin{align*}
    \widetilde{\rho}\left( \phi\psi \right) &= \lim_i \widetilde{\rho}\left( u_i\psi \right)\\
                                            &= \lim_i \rho\left( u_i \right)\rho\left( \psi \right)\\
                                            &= \widetilde{\rho}\left( \phi \right)\rho\left( \psi \right).
  \end{align*}
  In particular, we have $ \widetilde{\rho}\left( u_i\psi \right) = \rho\left( u_i \right) \widetilde{\rho}\left( \psi \right) $ for any $\psi\in B_{\infty}(X)$ for all $u_i\in C(X)$. Therefore, we have $ \widetilde{\rho}\left( u_i\psi \right) \rightarrow \widetilde{\rho}\left( \phi\psi \right) $ in WOT, whence
  \begin{align*}
    \widetilde{\rho}\left( \phi\psi \right) &= \widetilde{\rho}\left( \phi \right)\widetilde{\rho}\left( \psi \right)
  \end{align*}
  for any $\phi,\psi\in B_{\infty}(X)$.

  The linearity of measures gives that $ \widetilde{\rho} $ itself is linear. A similar bootstrapping argument gives $ \widetilde{\rho}\left( \phi \right)^{\ast} = \widetilde{\rho}\left( \overline{\phi} \right) $.

  Now, let $\Delta\subseteq X$ be Borel. Let $E\left( \Delta \right) = \widetilde{\rho}\left( \chi_{\Delta} \right)$. We will show that $E$ is a spectral measure.

  First, since $ \chi_{\Delta} $ is a hermitian idempotent in $B_{\infty}(X)$, it follows that $E\left( \Delta \right)$ is a projection by the classification of projections on a Hilbert space. Furthermore, $\chi_{\emptyset} = 0$ and $\chi_{X} = 1$, so that $E\left( \emptyset \right) = 0$ and $E\left( X \right) = 1$. Since
  \begin{align*}
    E\left( \Delta_1\cap \Delta_2 \right) &= \widetilde{\rho}\left( \chi_{\Delta_1\cap \Delta_2} \right)\\
                                          &= \widetilde{\rho}\left( \chi_{\Delta_1}\chi_{\Delta_2} \right)\\
                                          &= E\left( \Delta_1 \right)E\left( \Delta_2 \right).
  \end{align*}
  Finally, let $\set{\Delta_n : n\geq 0}$ be a pairwise disjoint sequence of Borel sets. Then, for fixed $h\in H$ and $\Lambda_n = \bigcup_{k=n+1}^{\infty}\Delta_k$, we have
  \begin{align*}
    \norm{E\left( \bigcup_{k=1}^{\infty}\Delta_k \right) - \sum_{k=1}^{n}E\left( \Delta_k \right)}^2 &= \iprod{E\left( \Lambda_n \right)h}{E\left( \Lambda_n \right)h}\\
               &= \iprod{E\left( \Lambda_n \right)h}{h}\\
               &= \int_{}^{} \chi_{\Lambda_n}\:d\mu_{h,h}\\
               &= \sum_{k=n+1}^{\infty}\mu_{h,h}\left( \Delta_k \right)\\
               &\xrightarrow{n\rightarrow\infty} 0.
  \end{align*}
  Thus, $E$ is a spectral measure.

  Finally, we must show that $\rho\left( u \right) = \int_{}^{} u\:dE$ for all $u\in C(X)$. In fact, we will show that $ \widetilde{\rho} \left( \phi \right) = \int_{}^{} \phi\:dE$ for all $\phi\in B_{\infty}(X)$. For a fixed $\phi\in B_{\infty}(X)$, we choose a Borel partition $\set{\Delta_1,\dots,\Delta_n}$ such that $\sup\set{\left\vert \phi(x) - \phi\left( x' \right) \right\vert : x,x'\in \Delta_k} < \ve$ for each $k$. Then,
  \begin{align*}
    \norm{\phi - \sum_{k=1}^{n}\phi\left( x_k \right)\chi_{\Delta_k}} &< \ve
  \end{align*}
  for any choices of $x_{k}\in \Delta_k$. Since $ \norm{\widetilde{\rho}} = 1 $, it follows that
  \begin{align*}
    \norm{ \widetilde{\rho}\left( \phi \right) - \sum_{k=1}^{n}\phi\left( x_k \right)E\left( \Delta_k \right) } &< \ve,
  \end{align*}
  so that $ \widetilde{\rho}\left( \phi \right) = \int_{}^{} \phi\:dE $ for any $\phi\in B_{\infty}\left( X \right)$.

  The uniqueness of $E$ follows from the fact that the measures $E_{g,h}$ are unique for each $g,h\in H$ as part of the Riesz Representation Theorem.
\end{theorem}
The primary use of the existence of spectral measures is to enable the Borel functional calculus for normal operators on Hilbert spaces.
\begin{theorem}
  Let $N$ be a normal operator on $H$. There is a unique spectral measure $E$ on the Borel subsets of $\sigma(N)$ such that:
  \begin{enumerate}[(a)]
    \item $N = \int_{}^{} z\:dE$;
    \item if $G$ is a nonempty open subset of $\sigma(N)$, then $E(G)\neq 0$;
    \item if $T\in B(H)$, then $TN = NT$ and $TN^{\ast} = N^{\ast}T$ if and only if $TE\left( \Delta \right) = E\left( \Delta \right)T$ for all Borel $\Delta\subseteq \sigma(N)$.
  \end{enumerate}
\end{theorem}
\begin{proof}
  Let $A = C^{\ast}\left( N \right)$. Then, $A$ is the closure of the polynomials in $N$ and $N^{\ast}$, so there is an isometric $\ast$-isomorphism $\rho\colon C\left( \sigma(N) \right)\rightarrow A$ given by the functional calculus, $\rho(u) = u(N)$. There is a unique spectral measure $E$ defined on the Borel subsets of $\sigma(N)$ such that $\rho(u) = \int_{}^{} u\:dE$ for all $u\in C\left( \sigma(N) \right)$. This gives (a).

  If $G$ is a nonempty open subset of $\sigma(N)$, Urysohn's Lemma yields a continuous function $u$ on $\sigma(N)$ such that $0\leq u(z)\leq \chi_G(z)$ for all $z\in \sigma(N)$. Since $E(G) = \widetilde{\rho}\left( \chi_G \right) \geq \rho(u)\neq 0$, it follows that $E(G)\neq 0$.

  Let $T\in B(H)$ be such that $TN = NT$ and $TN^{\ast} = N^{\ast}T$. Then, by the Stone--Weierstrass theorem, it follows that $Tu\left( N \right) = u\left( N \right)T$ for all $u\in C\left( \sigma(N) \right)$.

  Consider the set $\Omega$ defined as follows:
  \begin{align*}
    \Omega &= \set{\Delta\subseteq \sigma(N) : \Delta\text{ is a Borel subset and }AE\left( \Delta \right) = E\left( \Delta \right)A}.
  \end{align*}
  Notice that by the definition of the spectral measure, $\Omega$ is closed under finite intersections, countable unions (using SOT convergence), and contains $\sigma(N)$, so it is a $\sigma$-algebra by the Dynkin $\pi$-$\lambda$ theorem. If $G\subseteq \sigma(N)$ is open, there is a sequence $\left( u_n \right)_n$ of positive continuous functions on $\sigma(N)$ such that $u_n(z)\nearrow \chi_G(z)$ for all $z\in \sigma(N)$. Dominated convergence then gives
  \begin{align*}
    \iprod{AE(G)g}{h} &= \iprod{E(G)g}{A^{\ast}h}\\
                      &= E_{g,A^{\ast}h}\left( G \right)\\
                      &= \int_{}^{} \chi_G(z)\:dE_{g,A^{\ast}h}\\
                      &= \int_{}^{} \lim_{n\rightarrow\infty}u_n(z)\:dE_{g,A^{\ast}h}\\
                      &= \lim_{n\rightarrow\infty} \int_{}^{} u_n(z)\:dE_{g,A^{\ast}h}\\
                      &= \lim_{n\rightarrow\infty} \iprod{Au_n(N)g}{h}\\
                      &= \lim_{n\rightarrow\infty} \iprod{u_n(N)Ag}{h}\\
                      &= \iprod{E(G)Ag}{h}.
  \end{align*}
  Thus, $\Omega$ contains all the open subsets of $\sigma(N)$, so it must be the collection of Borel sets. The converse follows from bootstrapping by approximating with simple functions.
\end{proof}
\nocite{analysis_now,conway_functional_analysis}
\printbibliography
\end{document}
